\section{Conclusion}\label{sec:conclusion}

The audit of a small worked family turns a list of plausible claims about viability under incomplete observation into theorems with witnesses. The ordering of information classes is a theorem along three axes and strict along each; the two middle classes are incomparable; a no-repeat protocol with two instruments is alternation, and whether it costs anything depends on who keeps the record; decentralised unanimity factors statewise and loses viability through the allocation of a single floor vote; monitoring adequacy for the whole viable set is a graph condition with a closed-form count, with sixteen coarsest designs and none under a shared ledger; review timing is a max-plus hitting time whose identity must be written with the real slack; and the continuous fibre is certified with margins that are radii. Every number is regenerated from the primitives by exact code with independent implementations, and the ordering theorem is mechanised. The natural continuations are the three-instrument protocols, where alternation and the graph criterion both fail and new closed forms are needed, and the probabilistic refinement of the same audits.

\section*{Declarations}
\paragraph{Funding.} None.
\paragraph{Competing interests.} None.
\paragraph{Data availability.} No external data were used.
\paragraph{Code availability.} The exact core (\texttt{ws\_core.py}, \texttt{run\_all.py}), the table and figure generators, the verifier \texttt{paper11c\_worked\_systems\_audit\_v5\_verification.py} with its output, and the Lean modules are included in the submission package and available at \url{https://github.com/MIKEAA2020/general-sustainability}.
\paragraph{AI declaration.} Large language models assisted with drafting and iterative review of earlier editions; the author reviewed and edited all outputs and takes responsibility for the final work.

\begin{thebibliography}{99}
\footnotesize
\bibitem{Abaee2026calc} Abaee, A., 2026. An obstruction calculus for viability under incomplete observation. Figshare preprint, version 2. \url{https://doi.org/10.6084/m9.figshare.33716593.v2}.

\bibitem{Abaee2026cert} Abaee, A., 2026. Computational viability certification under incomplete observation: a continuous-to-finite bridge with exact certificates. Companion manuscript, \url{https://github.com/MIKEAA2020/general-sustainability}.

\bibitem{Abaee2026minimax} Abaee, A., 2026. Minimax dual certificates for the common-action obstruction: sparsity, polyhedral recovery, a benchmark law, and the finite review envelope. Companion manuscript, \url{https://github.com/MIKEAA2020/general-sustainability}.

\bibitem{Abaee2026prob} Abaee, A., 2026. Exactness is not the limitation: a linear criterion for blind survivable sets on the parameter cube and complete classifications for four, five and six parameters. Companion manuscript, \url{https://github.com/MIKEAA2020/general-sustainability}.

\bibitem{Althoff2020} Althoff, M., Bak, S., Bao, Z., Forets, M., Frehse, G., Freire, D., Kochdumper, N., Li, Y., Mitra, S., Ray, R., Schilling, C., Schupp, S., Wetzlinger, M., 2020. ARCH-COMP20 category report: continuous and hybrid systems with linear continuous dynamics. In: \emph{ARCH20, 7th International Workshop on Applied Verification of Continuous and Hybrid Systems}, EPiC Series in Computing 74, 16--48.

\bibitem{Aubin1991} Aubin, J.-P., 1991. \emph{Viability Theory}. Birkh\"auser, Boston.

\bibitem{Baccelli1992} Baccelli, F., Cohen, G., Olsder, G.J., Quadrat, J.-P., 1992. \emph{Synchronization and Linearity: An Algebra for Discrete Event Systems}. Wiley, Chichester.

\bibitem{Bernstein2002} Bernstein, D.S., Givan, R., Immerman, N., Zilberstein, S., 2002. The complexity of decentralized control of Markov decision processes. \emph{Mathematics of Operations Research} 27(4), 819--840.

\bibitem{BertsekasRhodes1971} Bertsekas, D.P., Rhodes, I.B., 1971. On the minimax reachability of target sets and target tubes. \emph{Automatica} 7(2), 233--247.

\bibitem{Blanchini1999} Blanchini, F., 1999. Set invariance in control. \emph{Automatica} 35(11), 1747--1767.

\bibitem{CassandrasLafortune2008} Cassandras, C.G., Lafortune, S., 2008. \emph{Introduction to Discrete Event Systems}, 2nd ed. Springer, New York.

\bibitem{CDHR2007} Chatterjee, K., Doyen, L., Henzinger, T.A., Raskin, J.-F., 2007. Algorithms for omega-regular games with imperfect information. \emph{Logical Methods in Computer Science} 3(3), paper 4.

\bibitem{CDH2010} Chatterjee, K., Doyen, L., Henzinger, T.A., 2010. Qualitative analysis of partially-observable Markov decision processes. In: \emph{Mathematical Foundations of Computer Science (MFCS 2010)}, Lecture Notes in Computer Science 6281, Springer, 258--269.

\bibitem{DanzerGrunbaumKlee1963} Danzer, L., Gr\"unbaum, B., Klee, V., 1963. Helly's theorem and its relatives. In: \emph{Convexity}, Proceedings of Symposia in Pure Mathematics 7, American Mathematical Society, 101--180.

\bibitem{Deffuant2007} Deffuant, G., Chapel, L., Martin, S., 2007. Approximating viability kernels with support vector machines. \emph{IEEE Transactions on Automatic Control} 52(5), 933--937.

\bibitem{Farkas1902} Farkas, J., 1902. Theorie der einfachen Ungleichungen. \emph{Journal f\"ur die reine und angewandte Mathematik} 124, 1--27.

\bibitem{Geretti2020} Geretti, L., Alexandre dit Sandretto, J., Althoff, M., Benet, L., Chapoutot, A., Chen, X., Collins, P., Forets, M., Freire, D., Immler, F., Kochdumper, N., Sanders, D.P., Schilling, C., 2020. ARCH-COMP20 category report: continuous and hybrid systems with nonlinear dynamics. In: \emph{ARCH20, 7th International Workshop on Applied Verification of Continuous and Hybrid Systems}, EPiC Series in Computing 74, 49--75.

\bibitem{Heidergott2006} Heidergott, B., Olsder, G.J., van der Woude, J., 2006. \emph{Max Plus at Work: Modeling and Analysis of Synchronized Systems}. Princeton University Press, Princeton.

\bibitem{Helly1923} Helly, E., 1923. \"Uber Mengen konvexer K\"orper mit gemeinschaftlichen Punkten. \emph{Jahresbericht der Deutschen Mathematiker-Vereinigung} 32, 175--176.

\bibitem{Ho1980} Ho, Y.-C., 1980. Team decision theory and information structures. \emph{Proceedings of the IEEE} 68(6), 644--654.

\bibitem{Kaelbling1998} Kaelbling, L.P., Littman, M.L., Cassandra, A.R., 1998. Planning and acting in partially observable stochastic domains. \emph{Artificial Intelligence} 101(1--2), 99--134.

\bibitem{LinWonham1988} Lin, F., Wonham, W.M., 1988. On observability of discrete-event systems. \emph{Information Sciences} 44(3), 173--198.

\bibitem{MouraUllrich2021} de Moura, L., Ullrich, S., 2021. The Lean 4 theorem prover and programming language. In: \emph{Automated Deduction -- CADE 28}, Lecture Notes in Computer Science 12699, Springer, 625--635.

\bibitem{Radner1962} Radner, R., 1962. Team decision problems. \emph{Annals of Mathematical Statistics} 33(3), 857--881.

\bibitem{Reif1984} Reif, J.H., 1984. The complexity of two-player games of incomplete information. \emph{Journal of Computer and System Sciences} 29(2), 274--301.

\bibitem{SaintPierre1994} Saint-Pierre, P., 1994. Approximation of the viability kernel. \emph{Applied Mathematics and Optimization} 29, 187--209.

\bibitem{SmallwoodSondik1973} Smallwood, R.D., Sondik, E.J., 1973. The optimal control of partially observable Markov processes over a finite horizon. \emph{Operations Research} 21(5), 1071--1088.

\bibitem{Stanley2012} Stanley, R.P., 2012. \emph{Enumerative Combinatorics, Volume 1}, 2nd ed. Cambridge University Press, Cambridge.

\bibitem{TsitsiklisAthans1985} Tsitsiklis, J.N., Athans, M., 1985. On the complexity of decentralized decision making and detection problems. \emph{IEEE Transactions on Automatic Control} 30(5), 440--446.

\bibitem{vonNeumann1928} von Neumann, J., 1928. Zur Theorie der Gesellschaftsspiele. \emph{Mathematische Annalen} 100, 295--320.

\bibitem{Witsenhausen1971} Witsenhausen, H.S., 1971. Separation of estimation and control for discrete time systems. \emph{Proceedings of the IEEE} 59(11), 1557--1566.
\end{thebibliography}

\end{document}
